%==============================================================================
% Sjabloon poster bachproef
%==============================================================================

\documentclass[a0,portrait]{hogent-poster}

% Info over de opleiding
\course{Bachelorproef}
\studyprogramme{Toegepaste Informatica}
\academicyear{2025-2026}
\institution{Hogeschool Gent, Valentin Vaerwyckweg 1, 9000 Gent}

% Info over de bachelorproef
\title{
    Vergelijking van Open Source LLM's voor\\
    Python Codegeneratie in Data Analyse en ETL
}

\author{Renaud Vermeiren}
\email{renaud.vermeiren@student.hogent.be}

\supervisor{Dhr. J. Decorte}
\cosupervisor{Dhr. D. Vanhaute}

\specialisation{AI \& Data Engineering}

\keywords{
    Large Language Models (LLMs),
    Code Generation,
    Software Reliability
}

\projectrepo{https://github.com/renaudVermeiren/Onderzoek}

\begin{document}
    
    \maketitle
    
    %==============================================================================
    % ABSTRACT
    %==============================================================================
    
    \begin{abstract}
        
        Deze bachelorproef onderzoekt hoe betrouwbaar zes populaire Large Language
        Models Python code genereren voor data analyse en ETL taken.
        Hiervoor werden 35 representatieve opdrachten uitgevoerd, goed voor
        630 geautomatiseerde code evaluaties op aspecten zoals syntaxis,
        beveiliging, onderhoudbaarheid en functionele correctheid.
        De resultaten tonen aan dat de meeste modellen technisch correcte en
        veilige code produceren, maar vaak tekortschieten op functionele
        correctheid. Vooral complexe ETL taken blijken een uitdaging,
        waardoor menselijke controle noodzakelijk blijft.
        
    \end{abstract}
    
    \begin{multicols}{2}
        
        %==============================================================================
        % INTRODUCTIE
        %==============================================================================
        
        \section{Introductie}
        
        Generatieve AI tools zoals ChatGPT, Copilot en Gemini worden steeds vaker
        gebruikt door studenten Toegepaste Informatica om Python code te genereren
        voor data analyse en ETL taken. Hoewel deze modellen snel bruikbare code
        kunnen produceren, is de kwaliteit en betrouwbaarheid van de gegenereerde
        oplossingen niet altijd gegarandeerd. Studenten beschikken vaak nog niet
        over voldoende ervaring om fouten, inefficiënties of onvolledige
        implementaties in deze code snel te herkennen. Daarom onderzoekt deze
        bachelorproef hoe zes open source Large Language Models presteren bij het
        genereren van Python code voor representatieve data analyse en ETL
        opdrachten.
        
        %==============================================================================
        % MODELLEN
        %==============================================================================
        
        
        
      \section{Onderzoeksopzet}
      
      Voor dit onderzoek werden zes open source Large Language Models geëvalueerd
      op hun vermogen om Python code te genereren voor data analyse en ETL taken.
      Elk model kreeg 35 representatieve opdrachten die aansluiten bij de
      leerinhoud van de afstudeerrichting Data Engineering en AI. Per opdracht
      werden drie oplossingen gegenereerd, wat resulteerde in 630 afzonderlijke
      code evaluaties.
      
      De gegenereerde code werd automatisch beoordeeld met verschillende
      evaluators:
      
      \begin{itemize}
          \item AST Parsing voor syntactische correctheid
          \item Ruff voor style controle
          \item Bandit voor beveiligingsanalyse
          \item Runtime tests voor uitvoerbaarheid
          \item Performance profiling voor performantie
          \item Radon voor onderhoudbaarheid
          \item Functionele testen voor correcte output
      \end{itemize}
        
        %==============================================================================
        % RESULTATEN
        %==============================================================================
        
        \section{Belangrijkste resultaten}
        
        \begin{itemize}
            \item 630 geautomatiseerde code evaluaties uitgevoerd
            \item 6 open source LLM's vergeleken
            \item Meer dan 95\% onderhoudbaarheid voor alle modellen
            \item Slechts 29.5\% tot 58.1\% functionele correctheid
            \item Grote verschillen tussen technische kwaliteit en correcte logica
        \end{itemize}
        
        \vspace{1cm}
        
        \begin{center}
            \captionsetup{type=figure}
            \includegraphics[width=0.95\linewidth]{01_evaluator_comparison}
            \captionof{figure}{
                Consistentie van de gegenereerde code per model en evaluator.
            }
        \end{center}
        
        %==============================================================================
        % CONCLUSIE
        %==============================================================================
        
        \section{Conclusie}
        
        De resultaten tonen aan dat de meeste modellen consistent presteren op
        syntaxis, beveiliging en onderhoudbaarheid. Voor uitvoerbaarheid en
        vooral functionele correctheid wordt echter een veel grotere variatie
        tussen samples waargenomen. Dit betekent dat dezelfde prompt niet altijd
        tot een even correcte oplossing leidt. Hoewel de onderzochte modellen
        vaak technisch correcte code produceren, blijft de betrouwbaarheid van
        de functionele uitkomst onvoldoende consistent om zonder menselijke
        controle te worden gebruikt voor data analyse en ETL taken. LLM's genereren vaak technisch correcte code,
        maar functionele correctheid blijft een belangrijke uitdaging.
        
      
        
        %==============================================================================
        % TOEKOMSTIG ONDERZOEK
        %==============================================================================
        
        \section{Toekomstig onderzoek}
        
        Dit onderzoek richtte zich op zes open source LLM's en een vaste set van
        data analyse en ETL taken. Toekomstig onderzoek kan nagaan welke invloed
        prompt engineering, andere temperature instellingen en geavanceerdere
        prompting technieken hebben op de kwaliteit van de gegenereerde code.
        Daarnaast kan het evaluatiekader worden uitgebreid met een grotere
        variatie aan ETL scenario's en recentere modellen.
        
    \end{multicols}
    
\end{document}